\begin{obeassessment}
\textbf{Kuis Bab 2}
\begin{enumerate}
    \item Apa perbedaan mendasar antara serangan pasif dan serangan aktif?
          Manakah yang lebih sulit dideteksi? Jelaskan alasannya.
    \item Dalam kondisi apa algoritma yang secara matematis kuat tetap dapat
          ditembus dengan serangan saluran samping?
    \item Diketahui cipherteks Caesar \code{KHOORZRUOG}. Serangan apa yang paling
          tepat untuk memecahkannya, dan bagaimana langkahnya? Jelaskan pula
          mengapa serangan \textit{chosen-plaintext} lebih kuat.
    \item Sebuah sistem mengirim perintah transfer tanpa nomor urut dan tanpa cap
          waktu. Kunci tetap rahasia dan cipher tidak dipecahkan, tetapi penyerang
          tetap dapat merugikan bank. Serangan apa yang dipakai, dan bagaimana
          penangkalannya?
    \item Jelaskan mengapa hasil pemulihan kunci melalui analisis daya (DPA)
          tidak dapat dicegah hanya dengan memilih algoritma yang lebih kuat.
\end{enumerate}
\end{obeassessment}
\begin{competencychecklist}
\begin{tabularx}{\linewidth}{|X|c|}
\hline
Indikator Kompetensi & Tercapai \\
\hline
Saya dapat menjelaskan prinsip Kerckhoffs dan mengapa kunci harus rahasia & [ ] \\
\hline
Saya dapat membedakan serangan berdasarkan informasi (COA, KPA, CPA, CCA) & [ ] \\
\hline
Saya dapat menyebutkan lima tujuan serangan & [ ] \\
\hline
Saya dapat membedakan serangan pasif dan aktif & [ ] \\
\hline
Saya dapat menjelaskan konsep keamanan komputasi & [ ] \\
\hline
Saya dapat menjelaskan empat jenis serangan saluran samping & [ ] \\
\hline
\end{tabularx}
\end{competencychecklist}
